\section{Every fixed Hankel concentration threshold}
\label{nh:section}

Fix a primitive finite-order Hecke character $\chi$ of
$K=\mathbb Q(\sqrt{-3})$, and retain the normalization
$c=2\pi/\sqrt{3N\mathfrak f_\chi}$ and the corrected theta map
$E_\chi$ from Section~\ref{sec:hecke-arithmetic-attempt}.
In positive variables, $E_\chi\psi(x)$ denotes
$(E_\chi\psi)(\log x)$.
Every fixed concentration threshold $0<\nu<1$ gives the same
all-zero trace limit as the deep selection. The estimate needed
for this assertion concerns exterior linear functionals in the
theta norm; it does not require a stronger bound for the whole
inverse Gram matrix.

On the raw space $L^2((0,\infty),du)$ write
\[
 P=P_\Lambda=1_{(0,\Lambda)},\quad
 B=H_cPH_c,\quad A=PBP,\quad F=PH_cP,\quad
 \Pi=\mathbf1_{[\nu,1]}(A).
\]
Thus $A=F^2$. Inverses of $A$ and $F$ below act only on
$\operatorname{Ran}\Pi$. Let $\mathcal R_\Lambda^\circ$ be this
range in a nonprincipal sector, and let
\[
 \mathcal R_\Lambda^\circ
 =\operatorname{Ran}\Pi\cap\ker I\cap\ker\operatorname{ev}_0
\]
in the principal sector. Denote its raw isometric inclusion by
$j_\Lambda$, and set
\[
 G_\Lambda=(E_\chi j_\Lambda)^*E_\chi j_\Lambda,\qquad
 Z_\Lambda=E_\chi j_\Lambda G_\Lambda^{-1/2},\qquad
 Q_\Lambda=Z_\Lambda Z_\Lambda^*.
\]
The Gram matrix is positive definite by the Mellin identity.
The fixed-threshold versions of
Proposition~\ref{prop:hankel-raw-estimates} follow from its proof
with $2$ replaced by $\nu^{-1}$. In particular the selected
profiles belong to the theta graph domain, are smooth up to zero
before their sharp endpoint, and admit the band extension
\begin{equation}\label{nh:extension}
 W_\Lambda=H_cPF^{-1}\Pi=BA^{-1}\Pi,\qquad
 PW_\Lambda=\Pi,\qquad \|F^{-1}\Pi\|\leq\nu^{-1/2}.
\end{equation}
The rank also has the full leading asymptotic. Indeed
\[
 |\operatorname{rank}\Pi-\Tr A|
 \leq\max\{\nu^{-1},(1-\nu)^{-1}\}\Tr(A-A^2)
 =O_{c,\nu}(\log(2+\Lambda)),
\]
by Proposition~\ref{prop:hankel-selection-size} and its variance proof.
Thus $\operatorname{rank}\Pi\sim2c\Lambda/\pi$.
All constants in this section may depend on the fixed character,
$c$, and $\nu$. Inner products are linear in their first argument.

\subsection{The two classical arithmetic estimates}

\begin{lemma}[Fixed-character reciprocal coefficients]
\label{nh:mertens}
Put
\[
 b_\chi(n)=\sum_{N\mathfrak a=n}
               \mu_K(\mathfrak a)\chi(\mathfrak a).
\]
There are constants $a_\chi,C_\chi>0$ such that
\begin{equation}\label{nh:mertens-bound}
 \left|\sum_{n\leq X}b_\chi(n)\right|
 \leq C_\chi X\exp\{-a_\chi\sqrt{\log(2+X)}\}
 \qquad(X\geq1).
\end{equation}
\end{lemma}

\begin{proof}
Apply the classical region of
\citet[Theorem~3.1, p.~1051]{thornerzaman2019} to a fixed abelian
class field through which $\chi$ factors. Its Dedekind zeta
function is the product of the associated primitive Hecke
functions. A zero of the chosen factor cannot be cancelled by
another factor away from the principal pole at one. The possible
exceptional zero is real, simple, and strictly less than one.
Since the character is fixed, decreasing a positive constant
excludes that one point. We obtain a zero-free region
\[
 \sigma\geq1-\frac{a_0}{\log(3+|t|)}
 \quad(s=\sigma+it).
\]
Nonprincipal Hecke factors are nonzero at one. In the principal
case $1/L_K(s,1)$ is holomorphic and zero there.

We record the reciprocal bound needed on the contour. For
$|t|\geq2$, put $\ell=\log(3+|t|)$ and choose $a>0$ sufficiently
small compared with $a_0$. On the disk centered at
$1+2a/\ell+it$ with radius $4a/\ell$, the analytic branch of
$\log L_K(s,\chi)$ exists. Convexity in this shrinking strip,
obtained from the functional equation and the line
$1+1/\ell$, gives
$\log|L_K(s,\chi)|=O_\chi(\log\ell)$ from above.
At the center the Euler logarithm is bounded in modulus by
$\log\zeta_K(1+2a/\ell)=O_\chi(\log\ell)$.
Borel--Carath\'eodory on the concentric radius $3a/\ell$ disk,
and the Euler product further to the right, yield
\[
 |L_K(\sigma+it,\chi)^{-1}|
       \ll_\chi(\log(3+|t|))^{C_\chi},
 \qquad 1-\frac a{\log(3+|t|)}\leq\sigma\leq2.
\]
Bounded heights follow by compactness and removability of the
principal reciprocal. The use of fixed constants here makes
no conductor-uniform assertion.

The quadratic-field ideal count gives $|b_\chi(n)|\leq d(n)$.
The divisor hyperbola argument gives, for $1\leq H\leq X$,
\[
 \sum_{X<n\leq X+H}d(n)\ll H\log(2X)+\sqrt X.
\]
Indeed, in every factor pair at least one factor is at most
$\sqrt{2X}$; sum the bounds $H/a+1$ for the other factor.
The same estimate holds for an interval immediately to the left
of $X$.

Take $X$ to be a half-integer,
$\kappa=1+1/\log X$, and $2\leq T\leq\sqrt X$.
The truncated Perron error is bounded by
\[
 \sum_{n\geq1}d(n)(X/n)^\kappa
       \min\{1,(T|\log(X/n)|)^{-1}\}.
\]
Outside $[X/2,2X]$, this is $O(X(\log X)^2/T)$.
Within that interval, first use the preceding short-interval
bound for $|n-X|\leq X/T$, then sum it over the doubled ranges
$2^jX/T<|n-X|\leq2^{j+1}X/T$.
Their total is $O(X(\log X)^2/T+\sqrt X)$. Hence
\[
 \sum_{n\leq X}b_\chi(n)
 =\frac1{2\pi i}\int_{\kappa-iT}^{\kappa+iT}
          \frac{X^s}{sL_K(s,\chi)}\,ds
 +O_\chi\left(\frac{X(\log X)^2}{T}+\sqrt X\right).
\]
Move the contour to $\sigma_0=1-a/\log(3+T)$, decreasing $a$
once more if necessary. The whole rectangle is zero-free and
there are no residues. The vertical integral is
\[
 O_\chi\left(
 X^{1-a/\log(3+T)}(\log(3+T))^{C_\chi+1}\right),
\]
and the horizontal integrals are
$O_\chi(X(\log(3+T))^{C_\chi}/T)$.
Set $T=\exp(b\sqrt{\log X})$ with a fixed $b>0$, and absorb
the logarithmic factors by decreasing the final positive
constant. This proves \eqref{nh:mertens-bound}.
Replacing an arbitrary $X$ by $\lfloor X\rfloor+1/2$ preserves
the summatory function and changes the bound by a fixed factor.
\end{proof}

\begin{lemma}[Mean square of the ideal discrepancy]
\label{nh:mean-square}
For
\[
 A_\chi(X)=\sum_{N\mathfrak a\leq X}\chi(\mathfrak a)
                        -\delta_\chi r_KX,
\]
one has
\begin{equation}\label{nh:mean-square-bound}
 \int_1^X|A_\chi(y)|^2\,dy\ll_\chi X^{3/2},
 \qquad
 \int_Y^\infty\frac{|A_\chi(y)|^2}{y^2}\,dy
                                      \ll_\chi Y^{-1/2}.
\end{equation}
\end{lemma}

\begin{proof}
We apply the established theorem of
\citet[Theorem~1, p.~62, equation~(57)]{chandrasekharannar1964}.
Its setup on pp.~43--44 permits complex coefficient sequences.
In its notation put
\[
 \varphi(s)=c^{-s}L_K(s,\chi),\qquad
 \psi(s)=\varepsilon_\chi c^{-s}L_K(s,\bar\chi),\qquad
 \Delta(s)=\Gamma(s),\qquad \delta=1.
\]
The functional equation is then
$\Delta(s)\varphi(s)=\Delta(1-s)\psi(1-s)$.
Thus its gamma parameter is $A=1$, and its two Dirichlet
scales are $\lambda_n=\mu_n=cn$. Absolute convergence and
the completed vertical-strip decay have the required form.
The only singularities are the usual finite set of poles.

Writing $a_\chi(n)=\sum_{N\mathfrak a=n}\chi(\mathfrak a)$,
we have $|a_\chi(n)|\leq d(n)$ and
\[
 \sum_{n\leq X}|a_\chi(n)|^2
 \leq\sum_{n\leq X}d_4(n)\ll X\log^3(2X).
\]
Here $d(n)^2\leq d_4(n)$ follows at each prime from
$(j+1)^2\leq\binom{j+3}{3}$.
The dual coefficients obey the same estimate.
Choose the theorem's parameters $\beta=1$, $\beta'=3$.
Its first branch explicitly includes
$2\beta-\delta-1/A=0$; the resulting parameter is
$\omega_0=1/4$. It gives
\[
 \int_1^X|E(y)|^2\,dy=C X^{3/2}+O(X\log^5X).
\]
Its residue convention gives
$E(y)=A_\chi(y/c)-L_K(0,\chi)$ away from the summation
jumps. The constant is bounded, and half weighting at the
discrete jumps changes no Lebesgue integral. Scaling back
and applying Cauchy--Schwarz proves the first assertion of
\eqref{nh:mean-square-bound}. Partial integration proves the
second assertion.
\end{proof}

\subsection{An exterior functional in the full theta norm}

\begin{lemma}[Exterior row estimate]\label{nh:row}
Let $q\in L^2((\Lambda,\infty),du)$ be extended by zero,
and put
\[
 k_q=F^{-1}\Pi PH_cq.
\]
Then $\|k_q\|_2\leq\nu^{-1/2}\|q\|_2$. For
$0<u\leq\Lambda/2$,
\begin{align}
 |k_q(u)|&\leq C\|q\|_2(1+\Lambda u)^{-1/4},
       \label{nh:pointwise}\\
 |u k_q'(u)|&\leq C\|q\|_2(1+\Lambda u)^{1/4}.
       \label{nh:derivative}
\end{align}
\end{lemma}

\begin{proof}
Let $J_u(v)=1_{(0,\Lambda)}(v)J_0(2c\sqrt{uv})$.
Since $F^{-1}\Pi=FA^{-1}\Pi$, evaluation of the finite
Hankel transform, followed by Cauchy--Schwarz, gives
\[
 |k_q(u)|^2
 \leq c^2\|q\|_2^2
       \langle J_u,A^{-2}\Pi(1-A)J_u\rangle
 \leq c^2\nu^{-2}\|q\|_2^2
       \langle J_u,(1-A)J_u\rangle.
\]
The use of $1-A$ is exact: $q$ is outside $P$ and
$P H_c(1-P)H_cP=P-A$.
Since $H_cPJ_u(v)=B_\Lambda(v,u)/c$,
\begin{equation}\label{nh:row-energy}
 |k_q(u)|^2\leq\nu^{-2}\|q\|_2^2
                 \int_\Lambda^\infty|B_\Lambda(v,u)|^2\,dv.
\end{equation}
Replacing $J_u$ by $u\partial_uJ_u$ gives the analogous
bound with $u\partial_uB_\Lambda(v,u)$.

The exact Lommel formula is
\[
 B_\Lambda(v,u)=\frac{c\sqrt\Lambda}{v-u}
 \left[
 \sqrt v J_1(2c\sqrt{\Lambda v})J_0(2c\sqrt{\Lambda u})
 -\sqrt u J_1(2c\sqrt{\Lambda u})J_0(2c\sqrt{\Lambda v})
 \right].
\]
For $v>\Lambda$ and $u\leq\Lambda/2$, the denominator is
at least $v/2$. The standard bounds for $J_0,J_1$ give
\[
 |B_\Lambda(v,u)|
 \leq C_c\Lambda^{1/4}v^{-3/4}(1+\Lambda u)^{-1/4}.
\]
Use
\[
 u\partial_uJ_0(2c\sqrt{\Lambda u})
       =-c\sqrt{\Lambda u}\,J_1(2c\sqrt{\Lambda u}),\qquad
 u\partial_u\{\sqrt u J_1(2c\sqrt{\Lambda u})\}
       =c\sqrt\Lambda\,uJ_0(2c\sqrt{\Lambda u})
\]
to obtain
\[
 |u\partial_uB_\Lambda(v,u)|
 \leq C_c\Lambda^{1/4}v^{-3/4}(1+\Lambda u)^{1/4}.
\]
The integral of $\Lambda^{1/2}v^{-3/2}$ over
$v>\Lambda$ is bounded. Substitution in
\eqref{nh:row-energy} proves both assertions.
\end{proof}

\begin{proposition}[Exterior functionals after Gram normalization]
\label{nh:functional}
For every $N>0$ and every $q$ supported in $(\Lambda,\infty)$,
\begin{equation}\label{nh:functional-bound}
 \sup_{0\ne\psi\in\mathcal R_\Lambda^\circ}
 \frac{|\langle W_\Lambda\psi,q\rangle|^2}
      {\|E_\chi\psi\|_2^2}
 \leq C_N\Lambda(\log\Lambda)^{-N}\|q\|_2^2.
\end{equation}
Equivalently, after relabelling $N$,
\begin{equation}\label{nh:normalized-extension}
 \|(1-P)W_\Lambda j_\Lambda G_\Lambda^{-1/2}\|
          \leq C_N\Lambda^{1/2}(\log\Lambda)^{-N}.
\end{equation}
\end{proposition}

\begin{proof}
Put $k=k_q$ from Lemma~\ref{nh:row}.
For $\psi\in\mathcal R_\Lambda^\circ$,
$\langle W_\Lambda\psi,q\rangle=\langle\psi,k\rangle$.
Write $b(\mathfrak a)=\mu_K(\mathfrak a)\chi(\mathfrak a)$.
Finite ideal inversion gives the exact pointwise identity
\[
 \psi(u)=\sum_{N\mathfrak a\leq\Lambda/u}
       \frac{b(\mathfrak a)}{\sqrt{N\mathfrak a\,u}}
                    (E_\chi\psi)(N\mathfrak a\,u).
\]
In the principal case the imposed mean condition makes the
corrected theta map equal to the ordinary theta sum here.
Define, for $0<v<\Lambda$,
\begin{equation}\label{nh:dual-representative}
 \beta_k(v)=\sqrt v\sum_{\mathfrak a}
       \frac{\overline{b(\mathfrak a)}}{N\mathfrak a}
                                      k(v/N\mathfrak a),
\end{equation}
and set $\beta_k=0$ above $\Lambda$.
We will prove
\begin{equation}\label{nh:dual-bound}
 \|\beta_k\|_{L^2(d^\times v)}^2
       \leq C_N\Lambda(\log\Lambda)^{-N}\|q\|_2^2.
\end{equation}

Separate the ideal of norm one. Its term is $\sqrt v\,k(v)$
and has squared $L^2(d^\times v)$ norm $\|k\|_2^2$.
For every remaining ideal, $v/N\mathfrak a\leq\Lambda/2$,
so \eqref{nh:pointwise}--\eqref{nh:derivative} apply.
Put $X=\Lambda v$. For large $X$ split the remaining sum at
$Y=X/(\log X)^D$, with $D$ a fixed number to be chosen.
The absolute ideal count is $O(y)$, and hence
\[
 \sum_{2\leq N\mathfrak a\leq Y}
     \frac{|k(v/N\mathfrak a)|}{N\mathfrak a}
 \ll\|q\|_2 X^{-1/4}Y^{1/4}
 =\|q\|_2(\log X)^{-D/4}.
\]
For the other terms apply Lemma~\ref{nh:mertens} and Abel
summation to $w(y)=y^{-1}k(v/y)$. Its derivative satisfies
\[
 |w'(y)|\leq C\|q\|_2y^{-2}(1+X/y)^{1/4}.
\]
The integral contribution is therefore bounded by
\[
 C\|q\|_2\int_Y^\infty
 e^{-a_\chi\sqrt{\log(2+y)}}(1+X/y)^{1/4}\frac{dy}{y},
\]
with a smaller boundary term. Since $X/Y=(\log X)^D$,
this is smaller than every fixed power of $(\log X)^{-1}$.
For bounded $X$, the same Abel argument gives $O(\|q\|_2)$.
Consequently the part of \eqref{nh:dual-representative}
from ideals of norm at least two obeys
\[
 |\beta_{k,\geq2}(v)|
 \leq C_N\|q\|_2\sqrt v\,
               [\log(2+\Lambda v)]^{-N}.
\]
Integrate its square against $dv/v$. The part
$v\leq\sqrt\Lambda$ is $O(\sqrt\Lambda\|q\|_2^2)$;
on the rest, $\log(2+\Lambda v)\asymp\log\Lambda$.
Adding the norm-one term and relabelling $N$ proves
\eqref{nh:dual-bound}.

We justify the dual identity without an exchange of
conditionally convergent sums. First integrate
$\langle\psi,k\rangle$ only over $u\geq\epsilon$.
Finite inversion and the substitution $v=N\mathfrak a\,u$
then give
\[
 \int_\epsilon^\Lambda\psi(u)\overline{k(u)}\,du
 =\int_0^\Lambda E_\chi\psi(v)\,
       \overline{\sqrt v
          \sum_{N\mathfrak a\leq v/\epsilon}
          \frac{\overline{b(\mathfrak a)}}{N\mathfrak a}
                         k(v/N\mathfrak a)}\,\frac{dv}{v}.
\]
The sum is finite. For fixed $\Lambda$, the truncated dual
functions converge to $\beta_k$ in $L^2(d^\times v)$.
Indeed, on $v\geq\delta>0$, Abel summation and
\eqref{nh:mertens-bound} make the reciprocal-coefficient
tail tend uniformly to zero. On $v<\delta$, smoothness of
$k$ at zero and the same Abel estimate give a uniform
$O_\Lambda(\sqrt v)$ bound. Its squared integral tends to
zero with $\delta$. The raw left side converges by
Cauchy--Schwarz. We have proved
\[
 \langle\psi,k\rangle_{L^2(du)}
       =\langle E_\chi\psi,\beta_k\rangle_{L^2(d^\times v)}.
\]
The estimate \eqref{nh:functional-bound} now follows from
\eqref{nh:dual-bound}. Duality gives
\eqref{nh:normalized-extension}.
\end{proof}

\subsection{The full smoothed theta tail}

\begin{lemma}[A raw Hilbert--Schmidt kernel]\label{nh:raw-tail}
Let $L=\log\Lambda$ and $g\in C_c^\infty(\mathbb R)$.
The operator
\[
 T_{\Lambda,g}=V(g)1_{t<-L}E_\chi(1-P)
\]
extends from the raw graph core to a Hilbert--Schmidt operator,
and
\begin{equation}\label{nh:raw-tail-bound}
 \|T_{\Lambda,g}\|_{\mathrm{HS}}
             \leq C_\chi\Lambda^{-1/2}
                             (\|g\|_2+\|g'\|_2).
\end{equation}
The same estimate holds for $g\in L^1(\mathbb R)\cap H^1(\mathbb R)$.
\end{lemma}

\begin{proof}
Use output coordinate $r=t+L$. For $u>\Lambda$, the raw
kernel is the Stieltjes integral
\[
 K_g(r,u)=u^{-1/2}\int_{\Lambda u}^\infty
        y^{-1/2}g(r+\log(y/(\Lambda u)))\,dA_\chi(y).
\]
For compact smooth $g$, the original sum is finite at each
fixed $(r,u)$. The principal density subtraction gives exactly
the discrepancy $A_\chi$ in this formula.
Integration by parts gives an endpoint
\[
 -\Lambda^{-1/2}u^{-1}A_\chi(\Lambda u)g(r)
\]
and an integral term
\[
 -\Lambda^{-1/2}u^{-1}
       \int_1^\infty A_\chi(\Lambda u z)z^{-3/2}
                    (g'-g/2)(r+\log z)\,dz.
\]
By Lemma~\ref{nh:mean-square}, the squared norm of the
endpoint is
\[
 \|g\|_2^2\int_{\Lambda^2}^\infty
                  |A_\chi(y)|^2y^{-2}\,dy
                       \ll_\chi\Lambda^{-1}\|g\|_2^2.
\]
Minkowski's inequality bounds the norm of the integral term by
\[
 \|g'-g/2\|_2
 \int_1^\infty z^{-3/2}
 \left(z\int_{z\Lambda^2}^\infty
                       |A_\chi(y)|^2y^{-2}\,dy\right)^{1/2}dz
 \ll_\chi\Lambda^{-1/2}\|g'-g/2\|_2
                         \int_1^\infty z^{-5/4}dz.
\]
This proves \eqref{nh:raw-tail-bound}.
The raw-kernel extension agrees with the theta graph
operator by approximation in $H_D^{3/8}$ on the exterior
half-line and the graph bound \eqref{eq:hecke-graph}.
Approximation of $g$ in $L^1\cap H^1$ gives the last assertion.
The $L^1$ condition ensures that convolution is bounded on
theta space; no such assertion is needed for an arbitrary
$H^1$ function.
\end{proof}

\begin{proposition}[Smoothed leakage at a fixed threshold]
\label{nh:leakage}
For every $N>0$ and every
$g\in C_c^\infty(\mathbb R)+\mathbb C\,(\tfrac12e^{-|t|})$,
\begin{equation}\label{nh:leakage-bound}
 \|V(g)(1-S_L)Z_\Lambda\|_{\mathrm{HS}}
                         \ll_{N,g}L^{-N}.
\end{equation}
\end{proposition}

\begin{proof}
The upper theta tail of a selected raw profile is zero.
In the principal case this uses $I\psi=0$.
By \eqref{nh:extension}, $H_cW_\Lambda\psi=F^{-1}\psi$
is supported in $(0,\Lambda)$; in the principal case
\[
 I(H_cW_\Lambda\psi)=W_\Lambda\psi(0)/c=\psi(0)/c=0.
\]
The graph-domain Poisson identity
$E_\chi H_c=\varepsilon_\chi J E_{\bar\chi}$ therefore
makes the lower theta tail of $W_\Lambda\psi$ zero.
Consequently
\[
 V(g)1_{t<-L}E_\chi\psi
       =-T_{\Lambda,g}(1-P)W_\Lambda\psi.
\]
Apply the Hilbert--Schmidt bound \eqref{nh:raw-tail-bound}
and the operator bound \eqref{nh:normalized-extension}.
Their factors $\Lambda^{-1/2}$ and $\Lambda^{1/2}$ cancel.
The stated Green kernel belongs to $L^1\cap H^1$,
so Lemma~\ref{nh:raw-tail} covers it as well.
\end{proof}

\subsection{All zero jets and the trace limit}

\begin{proposition}[Fixed zero packets]\label{nh:zero-packets}
Let $\mathcal V_{\mathcal F,L}$ be the truncated
exponential-polynomial space associated with a fixed finite
packet $\mathcal F$ of nontrivial zeros of $L_K(s,\chi)$,
including their full multiplicities. Then, for every $N>0$,
\begin{equation}\label{nh:zero-packet-bound}
 \|Q_\Lambda P_{\mathcal V_{\mathcal F,L}}\|_{\mathrm{HS}}^2
                                  \ll_{\mathcal F,N}L^{-N}.
\end{equation}
\end{proposition}

\begin{proof}
Fix $\rho=\beta+i\gamma$ in the packet and $0\leq j<m_\rho$.
The lower Mellin-tail kernel is
\[
 H_{\rho,j,\Lambda}(u)=
 \left.(\partial_s+L)^j\left[
 u^{s-1}\left(L_K(s,\chi)
       -\sum_{N\mathfrak a\leq\Lambda u}
                   \chi(\mathfrak a)(N\mathfrak a)^{-s}\right)
 +\delta_\chi\frac{r_K\Lambda^{1-s}}{1-s}
 \right]\right|_{s=\rho}.
\]
Proposition~\ref{prop:hecke-improper-mellin} justifies its
improper integral and all fixed derivatives. The kernel is in
raw $L^2$: below $\Lambda^{-1}$ its $L$-terms vanish at the
zero and its jets; above that point the endpoint discrepancy
and the Riesz estimate give the bounds used in
Proposition~\ref{prop:hecke-deep-zero-jets}.

We improve its exterior norm using
Lemma~\ref{nh:mean-square}. At $j=0$,
\[
 H_{\rho,0,\Lambda}(u)
 =-\Lambda^{-\rho}u^{-1}A_\chi(\Lambda u)
 +\rho u^{\rho-1}\int_{\Lambda u}^\infty
                             A_\chi(y)y^{-\rho-1}\,dy.
\]
The first term has squared exterior norm
\[
 \Lambda^{1-2\beta}\int_{\Lambda^2}^\infty
                         |A_\chi(y)|^2y^{-2}\,dy
                                      \ll\Lambda^{-2\beta}.
\]
For the integral use \eqref{eq:hecke-riesz}. One integration
by parts, valid for every fixed $\beta>0$, gives
\[
 \int_X^\infty A_\chi(y)y^{-\rho-1}\,dy
 =-\frac{\delta_\chi}{6\rho}X^{-\rho}
                         +O_{\chi,\rho}(X^{-\beta-1/4}).
\]
After multiplication by $\rho u^{\rho-1}$, the principal
term has squared exterior norm $O(\Lambda^{-2\beta-1})$.
The remainder has squared norm
\[
 O\left(\Lambda^{-2\beta-1/2}
                      \int_\Lambda^\infty u^{-5/2}\,du\right)
                         =O(\Lambda^{-2\beta-2}).
\]
For $j\geq1$, the operator $(\partial_s+L)^j$ kills the
endpoint factor $\Lambda^{-s}$ and the principal constant
term. Substitute $y=(\Lambda u)z$ in the Riesz remainder;
the remaining derivatives introduce only fixed powers of
$\log z$, integrable since $\beta+1/4>0$. In particular
\begin{equation}\label{nh:zero-tail-norm}
 \|1_{u>\Lambda}H_{\rho,j,\Lambda}\|_2^2
                                  \ll_{\rho,j}\Lambda^{-2\beta}.
\end{equation}

These kernels annihilate the fixed Hankel band, with value
zero at the origin in the principal case. Indeed, first use
inverse Hankel transforms of smooth compact band profiles;
Poisson makes their lower theta tail zero, and the disk
superposition proof of the improper Mellin formula applies
to these rapidly decreasing profiles. Such vectors are dense
in the band space. Evaluation at zero is continuous on that
space, so the principal value condition can be preserved.
The raw $L^2$ kernel extends the annihilation by continuity.
In particular its pairing with $W_\Lambda\psi$ is zero.
Its pairing with $\psi$ is therefore the negative of its
exterior pairing with $(1-P)W_\Lambda\psi$.

Apply Proposition~\ref{nh:functional} with
$q=1_{u>\Lambda}H_{\rho,j,\Lambda}$, taking conjugates
according to the Mellin pairing. The squared theta-dual norm
of the lower Mellin jet is
\[
 O_{\rho,j,N}\bigl(\Lambda^{1-2\beta}L^{-N}\bigr).
\]
The full Mellin jet vanishes at $\rho$, and the upper theta
tail is zero. Thus the same functional, with opposite sign,
is its window integral.

Put $\alpha=\beta-1/2$. For $\alpha<0$, use
$1_{[-L,L]}e^{(\alpha-i\gamma)(t+L)}(t+L)^j$.
The additional squared factor $\Lambda^{2\beta-1}$
cancels the displayed power. For $\alpha>0$, use the
right-endpoint normalization
$1_{[-L,L]}e^{(\alpha-i\gamma)(t-L)}(t-L)^j$.
Expand $(t-L)^j=((t+L)-2L)^j$; the additional factor is
$\Lambda^{1-2\beta}$, giving
$\Lambda^{2-4\beta}L^{O(j)-N}$.
On the central line, normalized polynomial jets introduce
only powers of $L$.
The finite endpoint Gram matrices have nonsingular limits;
opposite endpoints decouple, and the central polynomial blocks
have their nonsingular limiting Gram matrices.
Since $N$ is arbitrary, these fixed changes of basis prove
\eqref{nh:zero-packet-bound}.
\end{proof}

\begin{proposition}[Transfer to the uncompressed trace]
\label{nh:trace-transfer}
For every $h\in C_c^\infty(\mathbb R)$,
\[
 \Tr((Q_\Lambda-S_LQ_\Lambda S_L)V(h))\longrightarrow0.
\]
\end{proposition}

\begin{proof}
Let $\kappa(t)=\tfrac12e^{-|t|}$ and
$b=(1-\partial_t^2)h$. Then $h=b*\kappa$.
The isometry of $Z_\Lambda$ gives, for $g\in L^2$,
\[
 \|V(g)S_LZ_\Lambda\|_{\mathrm{HS}}
 \leq\|V(g)S_L\|_{\mathrm{HS}}
                         =\sqrt{2L}\,\|g\|_2.
\]
For $g$ in the class of Proposition~\ref{nh:leakage},
the same bound for $V(g)Z_\Lambda$ has an additional
$O_N(L^{-N})$ term.
Write $P_o=1-S_L$ and $b^*(t)=\overline{b(-t)}$.
Cyclicity of the finite-rank trace and
$Q-S_LQS_L=P_oQ+S_LQP_o$ give
\begin{align*}
 |\Tr((Q-S_LQS_L)V(h))|
 &\leq\|V(b^*)Z_\Lambda\|_{\mathrm{HS}}\,
                  \|V(\kappa)P_oZ_\Lambda\|_{\mathrm{HS}}\\
 &\quad+\|V(b^*)P_oZ_\Lambda\|_{\mathrm{HS}}\,
                  \|V(\kappa)S_LZ_\Lambda\|_{\mathrm{HS}}\\
 &\ll_{N,h}L^{1/2-N}+L^{-2N}.
\end{align*}
Here the two trace pairings are respectively
$Z_\Lambda^*V(b)V(\kappa)P_oZ_\Lambda$ and
$Z_\Lambda^*P_oV(b)V(\kappa)S_LZ_\Lambda$.
This proves the assertion without a square-root rank loss.
\end{proof}

\begin{theorem}[The fixed-threshold Hecke trace formula]
\label{nh:theorem}
For every fixed $0<\nu<1$, the compressed defect measures
associated with $\Pi=\mathbf1_{[\nu,1]}(A_\Lambda)$ converge
on all Schwartz frequency tests to the all-zero harmonic
measure $\mathfrak b_\chi$, with full zero multiplicities.
For every $h\in C_c^\infty(\mathbb R)$,
\begin{equation}\label{nh:trace-limit}
 2\log\Lambda\,h(0)-\Tr(Q_\Lambda V(h))
                              \longrightarrow\mathfrak B_\chi(h).
\end{equation}
The natural threshold $\nu=1/2$ is included.
In the principal full completion, the same rank-two polar
correction as in Proposition~\ref{prop:hecke-principal-graph}
must be added. No hypothesis on the positions or
multiplicities of the nontrivial zeros is used.
\end{theorem}

\begin{proof}
For a fixed finite zero packet put
$P_V=P_{\mathcal V_{\mathcal F,L}}$ and
$D_\Lambda=S_L-S_LQ_\Lambda S_L$.
Then $0\leq D_\Lambda\leq1$, $P_V\leq S_L$, and
\[
 \|(1-D_\Lambda^{1/2})P_V\|_{\mathrm{HS}}^2
 \leq\Tr(P_V(1-D_\Lambda)P_V)
                    =\|Q_\Lambda P_V\|_{\mathrm{HS}}^2=o(1)
\]
by Proposition~\ref{nh:zero-packets}.
For any nonnegative compactly supported continuous frequency test $\varphi$,
\[
 \Tr(D_\Lambda C_\varphi)
 =\|D_\Lambda^{1/2}C_\varphi^{1/2}\|_{\mathrm{HS}}^2
 \geq\|P_VD_\Lambda^{1/2}C_\varphi^{1/2}\|_{\mathrm{HS}}^2
 =\Tr(P_VC_\varphi)+o(1).
\]
Theorem~\ref{thm:finite-packet} identifies the last limit
with the harmonic measure of that packet. Increase the finite
packet and use positivity and the zero-count summability to
obtain the all-zero lower bound.

For large $\Lambda$, the deep raw selected space is contained
in the present fixed-threshold space, also after the principal
constraints. Their theta-image projections are therefore
nested, and
\[
 0\leq D_\Lambda^{(\nu)}\leq D_\Lambda^{\mathrm{deep}}.
\]
Theorem~\ref{thm:hi-local-limit} gives the matching compact
frequency upper bound. The uniform polynomial mass estimate
of Theorem~\ref{thm:hi-uniform-mass} passes through the same
order inequality. Equality of the compact limits, polarization,
and this uniform mass bound give convergence on every Schwartz
frequency test.
The compressed trace identity and
Proposition~\ref{nh:trace-transfer} prove
\eqref{nh:trace-limit}.

Finally take the two raw principal coordinate-repair vectors
already used for the deep cutoff in
Proposition~\ref{prop:hecke-principal-graph}.
They also belong to the present raw selected space. Let
$Y_\Lambda:\mathbb C^2\to L^2$ be their normalized theta
columns. The two graph maps are
\[
 X_\Lambda^{\mathrm{deep}}
       =(1-Q_\Lambda^{\mathrm{deep}})Y_\Lambda,\qquad
 X_\Lambda^{(\nu)}=(1-Q_\Lambda^{(\nu)})Y_\Lambda.
\]
Nesting gives
$\|X_\Lambda^{(\nu)}\|\leq
\|X_\Lambda^{\mathrm{deep}}\|\to0$.
The rank-two graph-projection formula consequently gives
trace-norm convergence of the full-principal correction to
the vertical $\mathbb C^2$ projection. Its integrated trace is
\[
 \int_{\mathbb R}h(t)(e^{-t/2}+e^{t/2})\,dt
\]
in the stated principal coordinate normalization. This proves
the full-sector assertion.
\end{proof}
